\documentclass[11pt]{article}
\usepackage[margin=1in]{geometry}
\usepackage[T1]{fontenc}
\usepackage[utf8]{inputenc}
\usepackage{lmodern}
\usepackage{microtype}
\usepackage[table]{xcolor}
\usepackage{amsmath,amssymb,mathtools}
\usepackage{booktabs,array,longtable}
\usepackage{hyperref}
\usepackage{enumitem}
\usepackage{fancyhdr}
\usepackage{titlesec}
\usepackage{tcolorbox}
\hypersetup{colorlinks=true, linkcolor=blue!50!black, urlcolor=blue!50!black}

\definecolor{TitleBlue}{HTML}{163A5F}
\definecolor{AccentTeal}{HTML}{2D6F73}
\definecolor{SoftBlue}{HTML}{EAF3F8}
\definecolor{SoftTeal}{HTML}{EDF8F6}
\definecolor{SoftGray}{HTML}{F5F7FA}

\pagestyle{fancy}
\fancyhf{}
\lhead{PINN for Two-Dimensional MHD Evolution}
\rhead{Implementation Note}
\cfoot{\thepage}
\setlength{\headheight}{14pt}

\titleformat{\section}{\Large\bfseries\color{TitleBlue}}{\thesection}{0.6em}{}
\titleformat{\subsection}{\large\bfseries\color{AccentTeal}}{\thesubsection}{0.6em}{}

\newtcolorbox{overviewbox}{colback=SoftBlue,colframe=TitleBlue,boxrule=0.8pt,arc=2mm,left=3mm,right=3mm,top=2mm,bottom=2mm}
\newtcolorbox{implbox}{colback=SoftTeal,colframe=AccentTeal,boxrule=0.8pt,arc=2mm,left=3mm,right=3mm,top=2mm,bottom=2mm}

\begin{document}

\begin{center}
    {\LARGE \bfseries \color{TitleBlue} Implementation Note: General Ideal-MHD PINN Residual Engine}\\[0.5em]
    {\large \color{AccentTeal} Physics-Informed Neural Network for Two-Dimensional MHD Evolution}
\end{center}

\vspace{0.5em}

\begin{overviewbox}
\textbf{Purpose.} This note explains what the new general residual engine adds relative to the reduced Alfv\'en PINN, and why it is included now as a scaffold rather than advertised as a finished full-PINN capability.
\end{overviewbox}

\section{What changed}
The earlier executable PINN in this project used a reduced shear-Alfv\'en subsystem. That was the right first step because it matched the first verification benchmark exactly and kept the training problem controllable.

The new file \texttt{pinn\_ideal\_mhd\_residuals\_general.m} broadens the residual definition to include the full primitive-form ideal-MHD system in two spatial dimensions with three vector components, often called 2.5D MHD.

\section{Residual groups now represented}
The scaffolded general engine includes residuals for:
\begin{enumerate}[leftmargin=1.8em]
    \item continuity,
    \item momentum in $x$,
    \item momentum in $y$,
    \item momentum in $z$,
    \item adiabatic pressure evolution,
    \item induction for $B_x$,
    \item induction for $B_y$,
    \item induction for $B_z$,
    \item magnetic divergence constraint.
\end{enumerate}

\section{Why primitive form was chosen here}
The network outputs primitive variables directly. Because of that, a primitive-form residual engine is easier to write and debug than a fully conservative residual engine derived from network-predicted conserved variables. The primitive form is therefore a reasonable intermediate step.

\begin{implbox}
\textbf{Caution.} Easier to implement does not mean easier to train. Primitive-form PINNs can be more sensitive to scaling and loss balancing, especially for strongly nonlinear coupled systems such as ideal MHD.
\end{implbox}

\section{Why the broader engine is not yet the main training path}
There are several technical reasons:
\begin{itemize}[leftmargin=1.8em]
    \item the general residual set is much stiffer than the reduced Alfv\'en system,
    \item density, pressure, and magnetic-field terms can produce losses of very different magnitudes,
    \item the training problem becomes much more sensitive to point sampling and normalization,
    \item stable and interpretable training usually requires additional infrastructure such as adaptive loss weights or staged curriculum training.
\end{itemize}

\section{How to extend it next}
The most sensible next extension is not to immediately throw all equations into one training run and hope for the best. The sensible extension is:
\begin{enumerate}[leftmargin=1.8em]
    \item normalize the fields and residuals carefully,
    \item test the general engine on a near-linear benchmark,
    \item add diagnostic plots for residual-group magnitudes during training,
    \item only then broaden to more nonlinear cases such as the magnetic pressure pulse.
\end{enumerate}

\begin{overviewbox}
\textbf{Bottom line.} The project now contains a scientifically honest bridge between the working reduced PINN and a broader general ideal-MHD PINN. That bridge is useful because it records the correct equations and derivative structure without pretending that the full training problem has already been solved.
\end{overviewbox}

\end{document}
